%% 2i2d-batterie-associations — trois branchements de quatre modules 12 V / 7 Ah.
%% Le branchement fixe le couple (U, Q) mais pas l'énergie E = U x Q, identique dans les trois cas.
%% Rail « + » en solideA, rail « - » en noir.
\documentclass[tikz,border=4pt]{standalone}
\usepackage{liaisons}
\begin{document}
\begin{tikzpicture}[x=1cm,y=1cm,
  film/.style={line width=0.8pt, line cap=round},          % fil côté -
  filp/.style={line width=0.8pt, line cap=round, draw=solideA}, % fil côté +
  titre/.style={font=\scriptsize\bfseries, anchor=south},
  val/.style={font=\scriptsize, anchor=north, inner sep=1pt},
  ener/.style={font=\scriptsize\bfseries, text=solideE, anchor=north, inner sep=2.5pt,
               draw=solideE, line width=0.7pt, rounded corners=2pt}]

%% ---- un module 12 V / 7 Ah, centré en (#1,#2) ----------------------------
\newcommand\modulebat[2]{%
  \draw[solideB, line width=0.7pt, fill=solideB!12] (#1-0.36,#2-0.30) rectangle (#1+0.36,#2+0.30);
  \node[font=\tiny] at (#1,#2+0.155) {12 V};
  \node[font=\tiny] at (#1,#2-0.145) {7 Ah};
  \draw[line width=0.7pt] (#1-0.32,#2) -- (#1-0.24,#2);                        % borne -
  \draw[filp] (#1+0.24,#2) -- (#1+0.32,#2)  (#1+0.28,#2-0.04) -- (#1+0.28,#2+0.04);}
%% ---- bornes du pack ------------------------------------------------------
\newcommand\bornemoins[2]{\fill (#1,#2) circle (1.1pt);
  \node[font=\scriptsize, anchor=south] at (#1,#2+0.04) {$-$};}
\newcommand\borneplus[2]{\fill[solideA] (#1,#2) circle (1.1pt);
  \node[font=\scriptsize, text=solideA, anchor=south] at (#1,#2+0.02) {$+$};}

%% ================= 1. PARALLÈLE ==========================================
\node[titre] at (0.55,3.45) {Parallèle};
\foreach \y in {2.77,2.03,1.29,0.55} {
  \modulebat{0.55}{\y}
  \draw[film] (0,\y) -- (0.19,\y);
  \draw[filp] (0.91,\y) -- (1.1,\y);}
\draw[film] (0,0.55) -- (0,3.00);
\draw[filp] (1.1,0.55) -- (1.1,3.00);
\bornemoins{0}{3.00}  \borneplus{1.1}{3.00}

%% ================= 2. SÉRIE-PARALLÈLE ====================================
\node[titre] at (2.5,3.45) {Série-parallèle};
\foreach \y in {2.05,1.25} {
  \modulebat{2.05}{\y}  \modulebat{2.95}{\y}
  \draw[film] (1.5,\y) -- (1.69,\y);
  \draw[line width=0.8pt] (2.41,\y) -- (2.59,\y);   % liaison série + vers -
  \draw[filp] (3.31,\y) -- (3.5,\y);}
\draw[film] (1.5,1.25) -- (1.5,3.00);
\draw[filp] (3.5,1.25) -- (3.5,3.00);
\bornemoins{1.5}{3.00}  \borneplus{3.5}{3.00}

%% ================= 3. SÉRIE ==============================================
\node[titre] at (5.95,3.45) {Série};
\foreach \x in {4.6,5.5,6.4,7.3} { \modulebat{\x}{1.66} }
\draw[film] (4.05,1.66) -- (4.24,1.66);
\draw[line width=0.8pt] (4.96,1.66) -- (5.14,1.66)  (5.86,1.66) -- (6.04,1.66)
                        (6.76,1.66) -- (6.94,1.66);
\draw[filp] (7.66,1.66) -- (7.85,1.66);
\draw[film] (4.05,1.66) -- (4.05,3.00);
\draw[filp] (7.85,1.66) -- (7.85,3.00);
\bornemoins{4.05}{3.00}  \borneplus{7.85}{3.00}

%% ================= tableau de synthèse ===================================
\rappel{\foreach \x/\u/\q in {0.55/12/28, 2.5/24/14, 5.95/48/7} {
  \node[val] at (\x,0.02) {$U = \u$ V};
  \node[val] at (\x,-0.33) {$Q = \q$ Ah};
  \node[ener] at (\x,-0.72) {$E = 336$ Wh};}

%% ---- accolade : la même énergie dans les trois cas ----------------------
\draw[solideE, line width=0.7pt] (0.55,-1.26) -- (5.95,-1.26);
\foreach \x in {0.55,2.5,5.95} { \draw[solideE, line width=0.7pt] (\x,-1.26) -- (\x,-1.18); }
\draw[solideE, line width=0.7pt] (3.25,-1.26) -- (3.25,-1.34);
\node[font=\scriptsize, text=solideE, anchor=north] at (3.25,-1.34)
  {$E = U \times Q$ : même énergie stockée dans les trois cas};}
\end{tikzpicture}
\end{document}
